\documentclass[11pt,a4paper]{article}

\usepackage[margin=1.05in]{geometry}
\usepackage{amsmath,amssymb,amsthm}
\usepackage{mathtools}
\usepackage{stmaryrd}
\usepackage{microtype}
\setlength{\emergencystretch}{2.5em}
\usepackage{enumitem}
\usepackage{booktabs}
\usepackage{array}
\usepackage{xcolor}
\usepackage[colorlinks=true,linkcolor=blue!45!black,urlcolor=blue!60!black]{hyperref}

\newtheorem{theorem}{Theorem}[section]
\newtheorem{proposition}[theorem]{Proposition}
\newtheorem{lemma}[theorem]{Lemma}
\newtheorem{corollary}[theorem]{Corollary}
\newtheorem{conjecture}[theorem]{Conjecture}
\theoremstyle{definition}
\newtheorem{definition}[theorem]{Definition}
\newtheorem{example}[theorem]{Example}
\theoremstyle{remark}
\newtheorem{remark}[theorem]{Remark}
\newtheorem{question}[theorem]{Question}
\newtheorem{claim}[theorem]{Claim}

\newcommand{\FF}{\mathbb{F}}
\newcommand{\NN}{\mathbb{N}}
\newcommand{\ZZ}{\mathbb{Z}}
\newcommand{\QQ}{\mathbb{Q}}
\newcommand{\CC}{\mathbb{C}}
\newcommand{\OO}{\mathcal{O}}
\newcommand{\cat}[1]{\mathbf{#1}}
\newcommand{\Set}{\cat{Set}}
\newcommand{\Vect}{\cat{Vect}}
\newcommand{\Coh}{\mathrm{Coh}}
\newcommand{\Bun}{\mathrm{Bun}}
\newcommand{\Rep}{\mathrm{Rep}}
\newcommand{\Perv}{\mathrm{Perv}}
\newcommand{\Gr}{\mathrm{Gr}}
\newcommand{\ext}[1]{\llbracket #1 \rrbracket}
\newcommand{\con}{\triangleleft}
\newcommand{\dual}[1]{{}^{L}\!#1}
\DeclareMathOperator{\Hom}{Hom}
\DeclareMathOperator{\End}{End}
\DeclareMathOperator{\GL}{GL}
\DeclareMathOperator{\SL}{SL}
\DeclareMathOperator{\SO}{SO}
\DeclareMathOperator{\Sp}{Sp}
\DeclareMathOperator{\Ind}{Ind}
\DeclareMathOperator{\Gal}{Gal}
\newcommand{\qbin}[2]{\genfrac{[}{]}{0pt}{}{#1}{#2}_{q}}
\newcommand{\Fl}{\mathrm{Fl}}
\newcommand{\Schub}{\mathfrak{S}}
\newcommand{\Groth}{\mathfrak{G}}

\definecolor{okgreen}{rgb}{0.0,0.45,0.15}
\definecolor{nored}{rgb}{0.65,0.0,0.0}
\definecolor{openblue}{rgb}{0.0,0.25,0.6}
\newcommand{\OK}{\textcolor{okgreen}{\textbf{verified}}}
\newcommand{\CL}{\textbf{classical}}
\newcommand{\NEW}{\textcolor{openblue}{\textbf{open}}}

\title{\bfseries The Hecke Algebra Was Already There\\[0.4em]
\large what \texttt{q-bool} and \texttt{lascoux} contain of the Langlands
programme, and the one chapter that closes the gap}
\author{Robin Langer \and Claude}
\date{24 September 2026}

\begin{document}
\maketitle

\begin{abstract}
Two repositories built for unrelated reasons --- \texttt{q-bool}, on containers
over a finite field, and \texttt{lascoux}, on the plactic monoid as a container
--- turn out to contain, unnamed, several of the standard objects of the
Langlands programme. The central observation is not an analogy: the operator
$S_i$ of \texttt{schubert.pdf}, defined there as pushforward along a forgetful
map of flags followed by pullback, \emph{is} the Iwahori--Hecke generator, and
the relations verified there are a presentation of $H_n(q)$. We re-derive this
from the geometry independently of \texttt{flag.py}
(\S\ref{sec:iwahori}); identify the refuted nilCoxeter conjecture's true target
as the Grothendieck polynomials of Lascoux--Sch\"utzenberger
(\S\ref{sec:groth}); locate the Langlands dual group in the Bernstein centre of
the \emph{affine} Hecke algebra, one level above where \texttt{q-bool}
currently works (\S\ref{sec:dual}); verify that the Hall constants of
\texttt{q-container.pdf}~\S9 are the Hall--Littlewood constants, which by
Macdonald's Chapter~V are the spherical Hecke algebra of $\GL_n$ over a local
field (\S\ref{sec:satake}); identify the two-parameter ceiling of
\texttt{q-container.pdf} with the loop direction, i.e.\ with Cherednik's double
affine Hecke algebra (\S\ref{sec:daha}); and show that
\texttt{q-dirichlet.pdf}~\S7 has already reached the automorphic side of
function-field Langlands (\S\ref{sec:autom}). For \texttt{lascoux} we record
three links, one of which --- $f^{\lambda}_{\mu\nu}(0) = c^{\lambda}_{\mu\nu}$
--- is a bridge between the two repositories rather than to Langlands, and is
verified here. Four computations, all independent of the theorems they check,
are in \texttt{code/verify.py}. \S\ref{sec:not} records what does \emph{not}
link, and \S\ref{sec:status} the status of every claim.
\end{abstract}

\section{The question}

\texttt{q-bool} was written to ask what containers become over $\FF_q$.
\texttt{lascoux} was written to ask whether Fomin's local rules are a container
isomorphism. Neither was written with the Langlands programme in view, and
neither mentions it.

The mishnah is therefore: nothing should connect. Container theory is a
fragment of dependent type theory; Langlands is arithmetic. The shared
vocabulary --- $q$, flags, Hall, Macdonald --- is the vocabulary of algebraic
combinatorics, which both subjects borrow from, and borrowing a word is not
sharing an object.

The objection is Proposition~4.1 of \texttt{schubert.pdf}. It records, as a
computational finding about an operator built from a fibration of flag
varieties, three relations:
\[
  T_i^2 = (q-1)T_i + qI,
  \qquad
  T_iT_jT_i = T_jT_iT_j \ (|i-j|=1),
  \qquad
  T_iT_j = T_jT_i \ (|i-j|\ge 2).
\]
That is not a family resemblance. It is a presentation, and the algebra it
presents has a name.

\section{Hecke operators, in one page}\label{sec:hecke}

For orientation, the object in its five standard guises. Only the first three
are needed below.

\paragraph{Sublattices.} A modular form of weight $k$ is a function on lattices
$\Lambda \subset \CC$, homogeneous of degree $-k$. For a prime $p$ there are
$p+1$ sublattices of index $p$, and no canonical choice among them, so one
sums:
\[
  (T_p F)(\Lambda) \;=\; \sum_{[\Lambda : \Lambda'] = p} F(\Lambda') .
\]
On $q$-expansions, $T_p f = \sum_n (a_{np} + p^{k-1}a_{n/p})q^n$; for an
eigenform normalised with $a_1 = 1$, the Fourier coefficients \emph{are} the
eigenvalues, which is why $L(s,f)$ has an Euler product.

\paragraph{Double cosets.} $T_p$ is the double coset
$[\Gamma \left(\begin{smallmatrix}1&0\\0&p\end{smallmatrix}\right)\Gamma]$.
Adelically, for a local field $F$ with ring of integers $\OO$ and residue field
$\FF_q$, the \emph{spherical} Hecke algebra is $C_c(K \backslash G / K)$ under
convolution with $K = G(\OO)$; the \emph{Iwahori} Hecke algebra takes instead
$K = $ the preimage of a Borel. The two differ, and the difference is exactly
the gap identified in \S\ref{sec:dual}.

\paragraph{Satake.} The spherical Hecke algebra is commutative, and
\[
  \mathcal{H}\bigl(G(F), G(\OO)\bigr) \;\cong\; \CC[X_*(T)]^{W}
  \;\cong\; \Rep\bigl(\dual{G}\bigr) ,
\]
the representation ring of the Langlands dual group. A Hecke eigenvalue is thus
not a number but a semisimple conjugacy class in $\dual{G}(\CC)$.

\paragraph{Trees.} For $G = \GL_2$, the quotient $G(F)/G(\OO)F^\times$ is the
vertex set of the $(q+1)$-regular Bruhat--Tits tree, and the Hecke operator is
its adjacency operator. Deligne's bound $|a_p| \le 2p^{(k-1)/2}$ is the optimal
spectral gap, which is where Ramanujan graphs come from.

\paragraph{Correspondences.} Geometrically, a Hecke operator modifies a
$G$-bundle at one point: the stack of $(E, E', x, \varphi)$ with $\varphi$ an
isomorphism away from $x$, with its two projections to $\Bun_G$. The type of
modification is a cocharacter of $G$, i.e.\ a weight of $\dual{G}$ --- Satake
again.

\medskip
\noindent
Each guise is the same refusal: an operation defined only up to a finite set of
choices at one place, and a sum instead of a choice.

\section{\texttt{schubert.pdf} constructs the Iwahori--Hecke
algebra}\label{sec:iwahori}

\begin{definition}\label{def:flag}
Let $\Fl_n(\FF_q)$ be the set of complete flags in $\FF_q^n$, $\pi_i$ the map
forgetting the $i$-th subspace, and
\[
  (S_i f)(F) \;=\; \sum_{\pi_i(F') = \pi_i(F)} f(F') .
\]
This is Definition~3.1 of \texttt{schubert.pdf}, verbatim.
\end{definition}

\begin{theorem}\label{thm:iwahori}
$S_i$ is convolution by the indicator function of the $i$-th minimal parabolic
$P_i$, hence $T_i := S_i - I$ is convolution by $\mathbf{1}_{Bs_iB}$; the $T_i$
generate the Iwahori--Hecke algebra $H_n(q)$ of type $A_{n-1}$, and
\[
  \End_{\GL_n(\FF_q)}\bigl(\Ind_B^G \mathbf{1}\bigr)
  \;=\; \End_{\GL_n(\FF_q)}\bigl(\CC[G/B]\bigr)
  \;\cong\; H_n(q) .
\]
\end{theorem}

\begin{proof}
Functions on $\Fl_n(\FF_q) = G/B$ with $G$ acting; $\pi_i$ is the projection
$G/B \to G/P_i$, so pushforward-then-pullback is convolution by
$\mathbf{1}_{P_i}$. The Bruhat decomposition of a minimal parabolic is
$P_i = B \sqcup Bs_iB$, so $\mathbf{1}_{P_i} = \mathbf{1}_B +
\mathbf{1}_{Bs_iB}$ and $T_i = S_i - I$ is convolution by
$\mathbf{1}_{Bs_iB}$. Since $|Bs_iB/B| = q$, the fibre of $\pi_i$ has
$q + 1 = |\mathbb{P}^1(\FF_q)|$ points and
$\mathbf{1}_{Bs_iB} * \mathbf{1}_{Bs_iB} = (q-1)\mathbf{1}_{Bs_iB} +
q\,\mathbf{1}_B$. The braid relations are the braid relations in $W = S_n$
lifted through the Bruhat decomposition. The last display is Iwahori's theorem
\cite{Iwahori}.
\end{proof}

\begin{remark}
This is what Proposition~4.1 of \texttt{schubert.pdf} records, one level of
abstraction down: there the relations are found by computing the matrix of
$S_i$ on cell classes and checking. The identification above says why they were
always going to hold, and that the cell rule quoted in Example~3.3 there,
\[
  S_i[w] =
  \begin{cases} [w] + [ws_i], & \ell(ws_i) > \ell(w),\\
                q[w] + q[ws_i], & \ell(ws_i) < \ell(w), \end{cases}
\]
is the standard $H_n(q)$ action on the Bruhat basis.
\end{remark}

\begin{proposition}\label{prop:v1}
\OK. \texttt{code/verify.py} rebuilds $\Fl_n(\FF_q)$ from subspace enumeration,
computes the relative position from the rank matrix of
Definition~\ref{def:flag}, checks $|\Fl_n| = [n]_q!$ and $|C_w| = q^{\ell(w)}$,
checks every fibre of $\pi_i$ has $q+1$ elements, checks that $S_i$ applied to
a cell indicator is constant on cells, and then verifies
$S_i^2 = (q+1)S_i$, the Hecke quadratic, both braid relations and
$T_i^2 \neq 0$, for $(n,q) \in \{(2,2),(2,3),(3,2),(3,3),(4,2)\}$ --- all $315$
flags and $24$ cells in the largest case. It shares no code with
\texttt{q-bool/code/flag.py} and does not assume the cell rule.
\end{proposition}

\section{The refutation relocated: Grothendieck, not
Schubert}\label{sec:groth}

\texttt{schubert.pdf} proves that no nilpotent divided difference exists on the
flag container over any finite field, and concludes that the nilCoxeter half of
its Conjecture is refuted. Both are correct. What the document does not say is
that the object it \emph{did} find has a Schubert-like basis of its own.

\begin{proposition}\label{prop:zero}
The $q \to 0$ degeneration of $H_n(q)$ is the $0$-Hecke algebra, generated by
idempotents $\pi_i$ with $\pi_i^2 = \pi_i$ --- Proposition~4.1(a) of
\texttt{schubert.pdf} is this statement for $S_i/(q+1)$. The $\pi_i$ are the
Demazure operators; they act on $K(G/B)$; and the basis they generate is the
\emph{Grothendieck polynomials} $\Groth_w$ of Lascoux--Sch\"utzenberger
\cite{LS82}.
\end{proposition}

So the flag container does not fail to produce a Schubert-type basis. It
produces the $K$-theoretic one. This is forced: counting $\FF_q$-points is a
$K$-theoretic invariant, and Schubert polynomials $\Schub_w$ live in the Chow
ring. The table in \S5 of \texttt{schubert.pdf} states exactly this dichotomy;
what is added here is that the left-hand column has a named polynomial basis,
and that the name is Lascoux --- the subject of the sibling repository.

\begin{proposition}\label{prop:nilhecke}
\texttt{schubert.pdf} Question~5.2 (grade the positions of the flag container
by the torus action; do the divided differences appear?) is answered by the
nilHecke ring of Kostant--Kumar \cite{KK86}, acting on $H^*_T(G/B)$, in which
$\partial_i$ and $\Schub_w$ live; Lusztig's graded affine Hecke algebra
\cite{Lusztig89} is the object interpolating between it and $H_n(q)$. The
question is not open.
\end{proposition}

\begin{remark}
The failure point $q = -1$ of \texttt{schubert.pdf} Theorem~4.2 is the
arithmetic signature of this dichotomy rather than an accident: $q = -1$ is
where a point count degenerates to an Euler characteristic.
\end{remark}

\section{Where the dual group is, and why \texttt{q-bool} cannot yet see
it}\label{sec:dual}

Theorem~\ref{thm:iwahori} gives the Iwahori-level algebra over a \emph{finite}
field. The Langlands dual group is not visible there.

\begin{proposition}\label{prop:centre}
The centre of the finite Hecke algebra $H_n(q)$ is the algebra of symmetric
polynomials in the Jucys--Murphy elements --- a shadow of $\Rep(\GL_n)$, but
finite-dimensional and carrying no lattice. By contrast, for the \emph{affine}
Hecke algebra $\widetilde{H}$ of $G(F)$, Bernstein's theorem gives
\[
  Z(\widetilde{H}) \;\cong\; \CC[X_*(T)]^{W} \;\cong\; \Rep(\dual{G}) .
\]
The Langlands dual group is the centre of the affine algebra, not of the finite
one.
\end{proposition}

\begin{remark}
\texttt{q-bool} is already at $\OO$ in one document and at $\FF_q$ in another,
and has not joined them. \texttt{q-container.pdf}~\S9 passes from
$\Vect_{\FF_q}$ to finite modules over a discrete valuation ring with residue
field $\FF_q$ --- the lattice side. \texttt{schubert.pdf} works with flags over
$\FF_q$ --- the flag side. The same local field underlies both. What is missing
is one construction.
\end{remark}

\begin{question}\label{q:affine}
\NEW. Define the \emph{affine flag container}: shapes the affine Weyl group
$\widetilde{W} = \ZZ^{n} \rtimes S_n$, positions $\FF_q^{\ell(w)}$, obtained by
replacing complete flags in $\FF_q^n$ by complete lattice chains in $\OO^n$.
Three things should then be checked, in this order.
\begin{enumerate}[nosep,label=(\roman*)]
\item The counting series $\sum_{w \in \widetilde{W}} X^{\ell(w)}$ is no longer
a polynomial. Bott's formula predicts a rational function; confirming which one
is the cheapest possible test that the construction is right.
\item The operators $S_i$, now indexed by $i = 0, 1, \dots, n-1$, satisfy the
affine braid relations together with the same quadratic.
\item The Bernstein centre of the resulting algebra is
$\CC[x_1^{\pm},\dots,x_n^{\pm}]^{S_n} = \Rep(\GL_n(\CC))$.
\end{enumerate}
Step (iii), if it goes through, is the Langlands dual group appearing inside a
container.
\end{question}

The $\GL_2$ case of the relevant spherical identity is verified in
\texttt{verify.py}, by enumerating sublattice chains in $\ZZ_p^2$ rather than
by quoting the Hecke ring:

\begin{proposition}\label{prop:v3}
\OK. For $p \in \{2,3,5\}$: there are $p+1$ sublattices of index $p$ and
$p^2+p+1$ of index $p^2$, of which exactly one has elementary divisors $(p,p)$
and $p^2+p$ have $(1,p^2)$; the number of chains
$\ZZ^2 \supset \Lambda' \supset \Lambda''$ with both steps of index $p$ is
$(p+1)^2$, distributed as multiplicity $p+1$ over the $(p,p)$ lattice and
multiplicity $1$ over each $(1,p^2)$ lattice. Hence
\[
  T_{(p,1)}^{\,2} \;=\; T_{(p^2,1)} \;+\; (q+1)\, T_{(p,p)} ,
\]
whose normalised Satake image is $(x+y)^2 = (x^2+y^2) + 2xy$ in
$\ZZ[x+y,(xy)^{\pm 1}] = \Rep(\GL_2)$.
\end{proposition}

\section{Hall--Littlewood is the spherical Hecke basis}\label{sec:satake}

This is the strongest link, and the cheapest to act on.

\begin{proposition}\label{prop:v2}
\OK. Let $g^{\lambda}_{\mu\nu}(p)$ count submodules $U \subseteq M_\lambda$ over
$\OO$ with $U \cong M_\nu$ and $M_\lambda/U \cong M_\mu$, and let
$f^{\lambda}_{\mu\nu}(t)$ be defined by
$P_\mu(x;t)P_\nu(x;t) = \sum_\lambda f^{\lambda}_{\mu\nu}(t) P_\lambda(x;t)$.
Then
\[
  g^{\lambda}_{\mu\nu}(p)
  \;=\; p^{\,n(\lambda) - n(\mu) - n(\nu)}\, f^{\lambda}_{\mu\nu}(p^{-1}) .
\]
Verified over $128$ instances, all $\lambda$ with $|\lambda| \le 4$, $p \in
\{2,3\}$: submodules enumerated by brute-force closure in $\ZZ/p^N$;
Hall--Littlewood polynomials computed symbolically from the Weyl-type formula in
four variables; neither side quoted from Macdonald.
\end{proposition}

That is Hall's theorem, and it is the content of \texttt{q-container.pdf}~\S9,
which is stated there as classical but unverified. What \S9 does not say is
what Chapter~V of the same book of Macdonald's says:

\begin{theorem}[Macdonald \cite{Macdonald}, Ch.~V]\label{thm:chV}
The spherical Hecke algebra $\mathcal{H}(\GL_n(F), \GL_n(\OO))$ of a local field
with residue field $\FF_q$ is isomorphic to the Hall algebra of finite
$\OO$-modules; its zonal spherical functions are the Hall--Littlewood
polynomials $P_\lambda(x;q^{-1})$; and composing with Satake gives
\[
  \text{Hall algebra of finite } \OO\text{-modules}
  \;\cong\; \mathcal{H}\bigl(\GL_n(F),\GL_n(\OO)\bigr)
  \;\cong\; \Rep\bigl(\dual{\GL_n}\bigr) .
\]
\end{theorem}

\begin{remark}
\texttt{q-container.pdf}~\S9 already writes the normalisation
$P_\lambda(x; q^{-1})$ --- Macdonald's Chapter~V normalisation --- and already
computes the Euler-form twist $q^{\langle M,N\rangle}$ relating the Hall algebra
to its twisted form. The book has one title because its two halves are one
subject, and the hinge is Chapter~V. Reading it against \S9 is the single
highest-yield next step in this document: it decides whether \texttt{q-bool}
contains a Langlands dual group or stops at Hall--Littlewood.
\end{remark}

\section{The two-parameter ceiling is the loop direction}\label{sec:daha}

\texttt{q-container.pdf}~\S\emph{Ceiling} observes that Macdonald's
$P_\lambda(x;q,t)$ has two parameters, that neither is the field size, that the
theory as built lands at $q_{\mathrm{Mac}} = 0$ (Hall--Littlewood), and that a
second direction is needed, with three candidates and no way to choose.

\begin{claim}\label{cl:daha}
The second direction is the loop. In Cherednik's double affine Hecke algebra
\cite{Cherednik}, of which Macdonald polynomials are the eigenfunctions:
$t$ is the Hecke parameter --- the field size, the $q+1$ of
Definition~\ref{def:flag}, the Iwahori level --- and $q$ is the parameter of the
translation lattice, i.e.\ loop rotation. The specialisation
$q_{\mathrm{Mac}} \to 0$ giving Hall--Littlewood is exactly the specialisation
that forgets the loop.
\end{claim}

\begin{remark}
The loop variable is already present in \texttt{q-container.pdf}, written $t$,
as the uniformiser of $\OO = \FF_q[[t]]$. If Claim~\ref{cl:daha} is right then
the document's own two letters are Macdonald's two parameters, and the
``second direction'' is not a new construction but a grading by $t$-adic degree
on a structure already defined. This is a prediction and is not verified here.
\end{remark}

\section{\texttt{q-dirichlet.pdf}~\S7 is the automorphic side}\label{sec:autom}

\texttt{q-dirichlet.pdf}~\S7 replaces the monoid of monic polynomials by the
monoid of effective divisors on a smooth projective curve $X/\FF_q$, obtains
$Z_X(u)$, and closes with a remark that categorifying the fibres --- replacing
divisors by coherent sheaves --- gives Kapranov's Hall algebra of $X$.

That remark is the door. Unramified automorphic forms for $\GL_n$ over the
function field $\FF_q(X)$ are functions on $\Bun_n(\FF_q)$; the Hecke operators
are modifications at closed points, which is \S\ref{sec:hecke} decategorified
from $D$-modules to point counts; and Kapranov \cite{Kapranov} makes the Hall
algebra of $\Coh(X)$ act by precisely these operators, recovering Eisenstein
series. This is the setting of Drinfeld's and Lafforgue's theorems.

\begin{remark}
The document's own observation that the Euler product is the statement that Day
convolution respects free generation is the right categorical account of why
automorphic $L$-functions factor over places; and the local factor at a closed
point $P$ is where the Satake parameter, a conjugacy class in $\dual{G}$, lives.
\end{remark}

\section{\texttt{lascoux}: three links, one of them internal}\label{sec:lascoux}

\paragraph{(a) Crystals are the geometric-Satake shadow of the dual group.}
Geometric Satake \cite{MV} gives
$\Perv_{G(\OO)}(\Gr_G) \cong \Rep(\dual{G})$ as tensor categories, and the
Mirkovi\'c--Vilonen cycles carry a canonical crystal structure isomorphic to
Kashiwara's \cite{BG,Kamnitzer}. \texttt{lascoux-container.pdf}~\S``crystals are
directed containers'' is therefore one functor from $\Rep(\dual{G})$. Note also
that $\Gr_G = G(F)/G(\OO)$ is the space on which the spherical Hecke algebra of
\S\ref{sec:satake} acts: crystals and Hall--Littlewood are two views of one
space.

\paragraph{(b) The plactic constants are the Hall constants at $t = 0$.}
This is a bridge between the two repositories that does not pass through
Langlands at all.

\begin{proposition}\label{prop:v4}
\OK. $f^{\lambda}_{\mu\nu}(0) = c^{\lambda}_{\mu\nu}$, the
Littlewood--Richardson coefficient. Verified over $64$ instances,
$|\lambda| \le 4$: the left side from the symbolic Hall--Littlewood expansion of
Proposition~\ref{prop:v2}, the right side by independent brute-force
enumeration of Littlewood--Richardson skew tableaux (rows weakly increasing,
columns strictly increasing, reverse reading word a lattice word), itself
sanity-checked against eight known coefficients.
\end{proposition}

\texttt{lascoux-container.pdf} asserts that the plactic Schur elements
$S_\lambda$ multiply by Littlewood--Richardson;
\texttt{q-container.pdf}~\S9 computes multiplication by
$g^{\lambda}_{\mu\nu}(q)$. Proposition~\ref{prop:v4} plus
Proposition~\ref{prop:v2} say these are the $t=0$ and generic-$t$ faces of one
structure. The plactic monoid is the degenerate fibre of the Hall algebra.

\paragraph{(c) The $q$-gap and the Macdonald ceiling are the same gap.}
The section of \texttt{lascoux-container.pdf} headed ``The gap: $q$, LLT, and
Hecke algebras'' states that there is no candidate route from a polynomial
functor to a quantum group.
LLT concerns Hecke algebras at roots of unity, and Ariki's theorem
\cite{Ariki} identifies their decomposition numbers with the canonical basis of
$U_v(\widehat{\mathfrak{sl}}_e)$ --- a quantum \emph{affine} algebra.

\begin{claim}\label{cl:samegap}
The missing direction in \texttt{lascoux} (LLT, quantum affine) and the missing
direction in \texttt{q-bool} (Macdonald's second parameter, Claim~\ref{cl:daha})
are the same direction: the loop. One gap, two repositories, two names.
\end{claim}

This is my claim, not a theorem. It is strongly suggested by the fact that both
resolutions are affine objects over the same base, and it predicts that
whichever repository first constructs the loop grading closes both gaps at once.

\section{What does not link}\label{sec:not}

Recorded so that the ledger is not only positive.

\begin{description}[leftmargin=1.6em,style=nextline,nosep]
\item[\texttt{finite-field-duality.pdf}.] The classification of nondegenerate
forms over $\FF_q$ is the classification of the groups on which the
non-self-dual Langlands duality acts ($\dual{\SO(2n+1)} = \Sp(2n)$), so the
input data coincide. But nothing in the document uses the root datum, and the
coincidence is of subject matter, not of construction. Weak.
\item[\texttt{error-correcting.pdf}, \texttt{data-control.pdf}.] No link found.
The data/control distinction is about variance in a container morphism; the
Hecke correspondence is not a container morphism in that sense, because both
legs move shapes. If there is a link it is through the flag container's algebra
of dispatch (\texttt{data-control.pdf}~\S4), which by
Theorem~\ref{thm:iwahori} is $H_n(q)$ --- but what that buys is not clear, and
is not claimed here.
\item[Arithmetic Langlands.] Nothing above reaches number fields. Everything is
over $\FF_q$ or a local/function field. The Galois side never appears.
\end{description}

\section{\emph{Mai nafka minah}}\label{sec:nafka}

What changes in practice, in order of cost.

\begin{enumerate}[nosep,label=\arabic*.]
\item Read Macdonald Chapter~V against \texttt{q-container.pdf}~\S9. Short
chapter, input already in your notation, decides
Theorem~\ref{thm:chV}'s applicability. Highest yield.
\item Amend \texttt{schubert.pdf}: name the algebra $H_n(q)$
(Theorem~\ref{thm:iwahori}), retitle the refutation around Grothendieck
polynomials (Proposition~\ref{prop:zero}), close Question~5.2 with
Kostant--Kumar (Proposition~\ref{prop:nilhecke}).
\item Add Proposition~\ref{prop:v4} to \texttt{lascoux}: the plactic constants
are the Hall constants at $t=0$. Cheap, verified, and it is the only statement
joining the two repositories.
\item Build the affine flag container (Question~\ref{q:affine}), test (i)
first. This is the only route to $\dual{G}$ from where \texttt{q-bool} stands.
\item Decide Claim~\ref{cl:daha} by computing whether the $t$-adic grading on
$\OO = \FF_q[[t]]$ induces Macdonald's $q$. If yes, Claim~\ref{cl:samegap}
follows and both gaps close.
\end{enumerate}

\section{Status of every claim}\label{sec:status}

\begin{center}
\footnotesize
\setlength{\tabcolsep}{2.5pt}
\begin{tabular}{clcl}
\toprule
\# & claim & where & status \\
\midrule
1 & $S_i$ is convolution by $\mathbf{1}_{P_i}$; $T_i$ generate $H_n(q)$
  & Thm.~\ref{thm:iwahori} & \CL{} (Iwahori); proof given \\
2 & relations hold, re-derived from the geometry
  & Prop.~\ref{prop:v1} & \OK{} --- $5/5$ cases, $315$ flags \\
3 & $q\to0$ gives Demazure, basis is $\Groth_w$
  & Prop.~\ref{prop:zero} & \CL{} (Lascoux--Sch.) \\
4 & Question~5.2 is Kostant--Kumar nilHecke
  & Prop.~\ref{prop:nilhecke} & \CL{}; the question is not open \\
5 & $\dual{G}$ is the Bernstein centre of the \emph{affine} algebra
  & Prop.~\ref{prop:centre} & \CL{} (Bernstein) \\
6 & $\GL_2$ spherical identity and its Satake image
  & Prop.~\ref{prop:v3} & \OK{} --- $p \in \{2,3,5\}$ \\
7 & Hall's theorem: $g^\lambda_{\mu\nu} \leftrightarrow$ Hall--Littlewood
  & Prop.~\ref{prop:v2} & \OK{} --- $128$ instances \\
8 & Hall algebra $\cong$ spherical Hecke $\cong \Rep(\dual{\GL_n})$
  & Thm.~\ref{thm:chV} & \CL{} (Macdonald Ch.~V); unverified \\
9 & Macdonald's second parameter is the loop
  & Cl.~\ref{cl:daha} & \CL{} (Cherednik); application \NEW{} \\
10 & \texttt{q-dirichlet}~\S7 is the $\GL_n$ automorphic side
  & \S\ref{sec:autom} & \CL{} (Kapranov); door left closed \\
11 & crystals $\leftrightarrow$ MV cycles $\leftrightarrow \Rep(\dual{G})$
  & \S\ref{sec:lascoux}(a) & \CL{} (Mirkovi\'c--Vilonen) \\
12 & $f^\lambda_{\mu\nu}(0) = c^\lambda_{\mu\nu}$: the repo bridge
  & Prop.~\ref{prop:v4} & \OK{} --- $64$ instances \\
13 & the LLT gap and the Macdonald gap are one gap
  & Cl.~\ref{cl:samegap} & \NEW{} --- my claim, unproved \\
14 & affine flag container has $\Rep(\dual{\GL_n})$ as its centre
  & Q.~\ref{q:affine} & \NEW{} --- not attempted \\
15 & \texttt{finite-field-duality} links to root data
  & \S\ref{sec:not} & weak; subject matter only \\
16 & anything reaches arithmetic Langlands
  & \S\ref{sec:not} & no \\
\bottomrule
\end{tabular}
\end{center}

\section{Related work}

\begin{thebibliography}{99}\small
\bibitem{Iwahori} N. Iwahori. \emph{On the structure of a Hecke ring of a
Chevalley group over a finite field}. J. Fac. Sci. Univ. Tokyo 10 (1964).
The theorem $\End_G(\Ind_B^G \mathbf{1}) \cong H_n(q)$, i.e.\
Theorem~\ref{thm:iwahori}.
\bibitem{LS82} A. Lascoux, M.-P. Sch\"utzenberger. \emph{Structure de Hopf de
l'anneau de cohomologie et de l'anneau de Grothendieck d'une vari\'et\'e de
drapeaux}. C. R. Acad. Sci. Paris 295 (1982). Grothendieck polynomials; the
basis the flag container actually produces.
\bibitem{KK86} B. Kostant, S. Kumar. \emph{The nil Hecke ring and cohomology of
$G/P$ for a Kac--Moody group}. Adv. Math. 62 (1986). Answers
\texttt{schubert.pdf} Question~5.2.
\bibitem{Lusztig89} G. Lusztig. \emph{Affine Hecke algebras and their graded
version}. J. Amer. Math. Soc. 2 (1989). The interpolating object.
\bibitem{Macdonald} I. G. Macdonald. \emph{Symmetric Functions and Hall
Polynomials}, 2nd ed., OUP 1995. Ch.~II--III is \texttt{q-container.pdf}~\S9;
Ch.~V is Theorem~\ref{thm:chV}.
\bibitem{Cherednik} I. Cherednik. \emph{Double Affine Hecke Algebras}. LMS
Lecture Notes 319, CUP 2005. Macdonald polynomials as DAHA eigenfunctions;
Claim~\ref{cl:daha}.
\bibitem{Kapranov} M. Kapranov. \emph{Eisenstein series and quantum affine
algebras}. J. Math. Sci. 84 (1997). The Hall algebra of a curve acting by
Hecke operators.
\bibitem{MV} I. Mirkovi\'c, K. Vilonen. \emph{Geometric Langlands duality and
representations of algebraic groups over commutative rings}. Ann. of Math. 166
(2007). Geometric Satake.
\bibitem{BG} A. Braverman, D. Gaitsgory. \emph{Crystals via the affine
Grassmannian}. Duke Math. J. 107 (2001).
\bibitem{Kamnitzer} J. Kamnitzer. \emph{Mirkovi\'c--Vilonen cycles and
polytopes}. Ann. of Math. 171 (2010).
\bibitem{Ariki} S. Ariki. \emph{On the decomposition numbers of the Hecke
algebra of $G(m,1,n)$}. J. Math. Kyoto Univ. 36 (1996). The LLT gap of
\texttt{lascoux}.
\bibitem{KW} A. Kapustin, E. Witten. \emph{Electric-magnetic duality and the
geometric Langlands program}. Commun. Number Theory Phys. 1 (2007). Background
for \S\ref{sec:hecke}; not used above.
\end{thebibliography}

\end{document}
